\textbf{(i) $\Rightarrow$ (ii)} If $S$ meets $\mathfrak{a}$, choose $s\in S\cap\mathfrak{a}$; both ideals are $A$. Otherwise $S$ avoids $r(\mathfrak{a})$, so the contraction is $\mathfrak{a}=(\mathfrak{a}:1)$ by Proposition 4.8.

\textbf{(ii) $\Rightarrow$ (iii)} Take $S=\{1,x,x^2,\ldots\}$. Its saturation is $\bigcup_{n\ge0}(\mathfrak{a}:x^n)$, and (ii) says that it equals one term of this increasing sequence.

\textbf{(iii) $\Rightarrow$ (i)} Pass to $A/\mathfrak{a}$ and take $\mathfrak{a}=0$. Suppose $xy=0$ with $y\ne0$, and choose $n$ such that $(0:x^n)=(0:x^{n+1})$. If $a\in(x^n)\cap(y)$, write $a=bx^n=cy$. Then $bx^{n+1}=ax=0$, so $bx^n=0$ and $a=0$. Irreducibility of zero and $(y)\ne0$ give $(x^n)=0$. Thus every zero-divisor is nilpotent.
